%% kapa: exam=1 round=24 year=2013 period=1 form=A subject=civil_law subject_ko=민법 seq=1 number=1
%% source: Q-Net 감정평가사 기출문제 (cst00309) artlSeq=5073802 file=extracted/1008/5073802/2013년 제24회 감정평가사 1차시험 기출문제/1교시(A형).hwp
%% conversion: kordoc 4.24.51 -> tools/gen1.py (auto); review=proofread-sonnet; flags=none
\begin{question}{1}
신의성실의 원칙에 관한 설명으로 옳지 않은 것은? (다툼이 있으면 판례에 의함)
\begin{choices}
  \item 당사자의 주장이 없으면 법원은 직권으로 신의칙의 위반 여부를 판단할 수 없다.
  \item 특정채무를 보증한 경우 채권자의 권리행사가 신의칙에 비추어 용납할 수 없는 때에는 극히 예외적으로 보증인의 책임을 제한할 수 있다.
  \item 신의칙은 사법(私法) 전반에 적용되는 일반원칙이다.
  \item 신의칙은 법률행위해석의 기준이 된다.
  \item 부동산거래에서 거래 상대방이 일정한 사정에 관한 고지를 받았더라면 그 거래를 하지 않았을 것이 경험칙상 명백한 경우 그 사정을 고지할 신의칙상 의무가 인정된다.
\end{choices}
\end{question}
\answer{1}
